\documentclass[11pt]{article}
\usepackage[T1]{fontenc}
\usepackage[utf8]{inputenc}
\usepackage{lmodern,microtype}
\usepackage[margin=1.08in]{geometry}
\usepackage{amsmath,amssymb,amsthm,mathtools}
\usepackage{xcolor,hyperref}
\hypersetup{colorlinks=true,linkcolor=blue!55!black,citecolor=blue!55!black,urlcolor=blue!60!black}
\setlength{\parindent}{0pt}
\setlength{\parskip}{0.48em}
\setlength{\emergencystretch}{2em}
\newtheorem{theorem}{Theorem}
\newtheorem{lemma}[theorem]{Lemma}
\newtheorem{corollary}[theorem]{Corollary}
\theoremstyle{remark}
\newtheorem{remark}[theorem]{Remark}
\newcommand{\Area}{\operatorname{Area}}
\newcommand{\Ric}{\operatorname{Ric}}

\title{A Positively Curved Counterexample to a CMC Index--Area Bound}
\author{[Author Name]}
\date{}

\begin{document}
\maketitle

\begin{abstract}
\textbf{Relation to the problem list.}
This manuscript addresses Question~3 (L3), ``CMC index versus area and
genus.''  Relative to the hypotheses printed in the list, it gives a full
counterexample: one fixed complete ambient three-manifold has strictly
positive sectional curvature and violates every possible constant $C(N)$.
The example is noncompact and its curvature decays to zero, so it does not
disprove variants requiring a closed ambient manifold or a uniform positive
curvature lower bound.

We construct a complete rotationally symmetric metric of strictly positive
sectional curvature on $\mathbb R^3$ containing embedded constant-mean-
curvature spheres whose areas diverge while their strong Morse index is one
and their nullity is zero.  This disproves an index--area estimate proposed
for branched CMC immersions when the ambient three-manifold is assumed only
to have pointwise positive sectional curvature.  The example isolates the
missing global hypothesis: the ambient curvature is positive everywhere but
decays to zero along the end.
\end{abstract}

\section{Introduction}

For a two-sided CMC immersion $u\colon\Sigma^2\to N^3$, let
\[
 L=-\Delta_\Sigma-\bigl(|A|^2+\Ric_N(\nu,\nu)\bigr)
\]
be the strong Jacobi operator.  Denote its index and nullity by
$i_{A_h}(u)$ and $n_{A_h}(u)$.  Seemungal and Sharp proposed that, when
$N$ has strictly positive sectional curvature, there should exist a
constant $C(N)>0$ such that
\begin{equation}\label{eq:proposed}
 C(N)\left((1+h^2)\Area(\Sigma)+g(\Sigma)\right)
 \le i_{A_h}(u)+n_{A_h}(u)
\end{equation}
for every branched CMC-$h$ immersion \cite{SeemungalSharp}.  We show that
pointwise positivity of the ambient sectional curvature does not suffice.

\begin{theorem}\label{thm:main}
There is a fixed complete orientable Riemannian three-manifold $(N,g)$ with
strictly positive sectional curvature and a sequence of embedded CMC
spheres $\Sigma_j\subset N$ such that
\[
 \Area(\Sigma_j)\longrightarrow\infty,
 \qquad
 i_{A_{h_j}}(\Sigma_j)=1,
 \qquad
 n_{A_{h_j}}(\Sigma_j)=0.
\]
Consequently, no positive constant $C(N)$ can make
\eqref{eq:proposed} hold on this ambient manifold.
\end{theorem}

\section{A complete positively curved model}

Fix $a\in(0,1)$ and define
\begin{equation}\label{eq:warping}
 f(r)=r(1+r^2)^{(a-1)/2},\qquad r\ge0.
\end{equation}
On $N=\mathbb R^3$, consider the rotational metric
\begin{equation}\label{eq:metric}
 g=dr^2+f(r)^2g_{\mathbb S^2}.
\end{equation}

\begin{lemma}\label{lem:geometry}
The metric \eqref{eq:metric} extends smoothly across $r=0$, is complete,
and has strictly positive sectional curvature everywhere.
\end{lemma}

\begin{proof}
The function $f$ is the restriction of a smooth odd function and
\[
 f(r)=r+\frac{a-1}{2}r^3+O(r^5).
\]
Thus the standard pole conditions $f(0)=0$ and $f'(0)=1$ hold, and
\eqref{eq:metric} extends smoothly across the origin.  Since the radial
coefficient is identically one on $[0,\infty)$, the metric is complete.

Direct differentiation gives
\begin{align}
 f'(r)&=(1+r^2)^{(a-3)/2}(1+ar^2)>0,\label{eq:fp}\\
 f''(r)&=(a-1)r(1+r^2)^{(a-5)/2}(3+ar^2)<0
 \quad(r>0).\label{eq:fpp}
\end{align}
The radial and tangential sectional curvatures of \eqref{eq:metric} are
\[
 K_{\rm rad}=-\frac{f''}{f},
 \qquad
 K_{\rm tan}=\frac{1-(f')^2}{f^2}.
\]
Equation \eqref{eq:fpp} gives $K_{\rm rad}>0$.  Moreover $f'$ decreases
strictly from $1$, so $0<f'(r)<1$ for $r>0$ and $K_{\rm tan}>0$.
At the pole both curvatures extend continuously to
\[
 K_{\rm rad}(0)=K_{\rm tan}(0)=3(1-a)>0.
\]
\end{proof}

\section{Coordinate spheres and their Jacobi spectra}

For $R>0$, let $\Sigma_R=\{r=R\}$.  Its induced metric is
$f(R)^2g_{\mathbb S^2}$, its genus is zero, and
\begin{equation}\label{eq:area-h}
 \Area(\Sigma_R)=4\pi f(R)^2,
 \qquad
 h_R=2\frac{f'(R)}{f(R)},
\end{equation}
up to the irrelevant choice of orientation.  Thus every $\Sigma_R$ is an
embedded CMC sphere.  Since $f(R)\sim R^a$, its area tends to infinity.

\begin{lemma}\label{lem:spectrum}
For all sufficiently large $R$,
\[
 i_{A_{h_R}}(\Sigma_R)=1,
 \qquad
 n_{A_{h_R}}(\Sigma_R)=0.
\]
\end{lemma}

\begin{proof}
For a coordinate sphere,
\[
 |A|^2=2\left(\frac{f'}f\right)^2,
 \qquad
 \Ric_N(\nu,\nu)=-2\frac{f''}f.
\]
The Jacobi potential is therefore constant and equals
\[
 q_R=\frac{2\bigl((f')^2-ff''\bigr)}{f^2}.
\]
The eigenvalues of $-\Delta_{\Sigma_R}$ are
$\ell(\ell+1)/f(R)^2$, with multiplicity $2\ell+1$.  Hence the complete
Jacobi spectrum is
\begin{equation}\label{eq:spectrum}
 \lambda_\ell(L_R)
 =\frac{\ell(\ell+1)-S(R)}{f(R)^2},
 \qquad
 S(R)=2\bigl(f'(R)^2-f(R)f''(R)\bigr).
\end{equation}
Using $f(R)\sim R^a$, $f'(R)\sim aR^{a-1}$, and
$f''(R)\sim a(a-1)R^{a-2}$ gives
\[
 S(R)\sim 2aR^{2a-2}\longrightarrow0^+.
\]
For all large $R$ one has $0<S(R)<2$.  In \eqref{eq:spectrum}, the
$\ell=0$ eigenvalue is then negative and every $\ell\ge1$ eigenvalue is
positive.  This proves the asserted index and nullity.
\end{proof}

\section{Failure of the estimate}

\begin{proof}[Proof of Theorem~\ref{thm:main}]
Choose any sequence $R_j\to\infty$ for which Lemma~\ref{lem:spectrum}
applies.  Equations \eqref{eq:area-h} give
\[
 (1+h_{R_j}^2)\Area(\Sigma_{R_j})
 =4\pi f(R_j)^2+16\pi f'(R_j)^2
 \longrightarrow\infty,
\]
whereas
\[
 i_{A_{h_{R_j}}}(\Sigma_{R_j})+n_{A_{h_{R_j}}}(\Sigma_{R_j})=1.
\]
The ambient manifold is fixed, so no positive constant depending only on
$N$ can satisfy \eqref{eq:proposed} for this sequence.
\end{proof}

\begin{remark}
The curvature of \eqref{eq:metric} is strictly positive but tends to zero
at infinity.  Thus the construction does not address variants in which the
ambient manifold is closed or has a uniform positive curvature lower bound.
\end{remark}

\begin{thebibliography}{9}

\bibitem{SeemungalSharp}
K.~Seemungal and B.~Sharp,
\emph{Index estimates for constant mean curvature surfaces in three-manifolds by energy comparison},
\href{https://arxiv.org/abs/2411.02932}{arXiv:2411.02932}.

\bibitem{ONeill}
B.~O'Neill,
\emph{Semi-Riemannian Geometry with Applications to Relativity},
Academic Press, New York, 1983.

\end{thebibliography}
\end{document}
